\documentclass[11pt, a4paper]{article}

\usepackage[utf8]{inputenc}
\usepackage[a4paper, total={6in, 8in}]{geometry}
\raggedbottom

\usepackage{mlmodern}
\usepackage{csquotes}
\usepackage{parskip}
\usepackage[italian]{babel}
\usepackage{booktabs}
\usepackage{multicol}
\usepackage{graphicx}
\usepackage{fancyhdr}
\usepackage{float}
\usepackage{changepage}
\usepackage{fancyvrb}
\usepackage[final]{pdfpages}
\usepackage{placeins}
\usepackage{hyperref}
\usepackage[all]{hypcap}
\usepackage{tcolorbox}
\usepackage{xurl}
\usepackage{xcolor}
\usepackage{tikz}
\usepackage{biblatex}
\usepackage{array}
\usepackage{pdfpages}
\usepackage{subfig}
\usepackage{multirow}
\usepackage{pgfplots}
\usepackage{listings}
\usepackage{subfig}
\usepackage{caption}
\usepackage{pdfpages}
\usepackage{placeins}
\usepgfplotslibrary{groupplots}
\pgfplotsset{compat=1.18}


\addbibresource{bibliography.bib}

\hypersetup{
	colorlinks=false,
	citecolor=black,
	filecolor=black,
	linkcolor=black,
	urlcolor=black,
	allbordercolors=white
}


\definecolor{codegreen}{rgb}{0,0.6,0}
\definecolor{codegray}{rgb}{0.5,0.5,0.5}
\definecolor{codepurple}{rgb}{0.58,0,0.82}
\definecolor{backcolour}{rgb}{0.95,0.95,0.92}

\lstdefinestyle{mystyle}{
	backgroundcolor=\color{backcolour},   
	commentstyle=\color{codegreen},
	numberstyle=\tiny\color{codegray},
	stringstyle=\color{codepurple},
	basicstyle=\ttfamily\footnotesize,
	breakatwhitespace=false,         
	breaklines=true,                 
	captionpos=b,                    
	keepspaces=true,                   
	numbersep=5pt,                  
	showspaces=false,                
	showstringspaces=false,
	showtabs=true,                  
	tabsize=1
}

\lstset{style=mystyle}

\setlength{\headheight}{14pt}

\begin{document}
	\begin{center}
	{\LARGE \textbf{Università degli Studi di Milano Bicocca}} \\
	\vspace{0.2cm}
	{\Large {Dipartimento di Informatica}} \\ 
	\vspace{1cm}
	
	\includegraphics[width=4cm]{figures/unimib-logo.png} \\
	\vspace{0.4cm}
	
	{\huge \textbf{Metodi del calcolo scientifico}} \\ 
	\Large{Compressione di immagini attraverso la DCT}
	\vspace{3cm}
\end{center}
\vspace{2cm}

\noindent \large{Membri del gruppo:} 

\begin{tabular}{@{}lll}
	\textbf{Canesi Gabriele} & 851637 & \href{mailto:g.canesi1@campus.unimib.it}{\texttt{g.canesi1@campus.unimib.it}}\\
	\textbf{Fiorenza Gioele} & 851631 & \href{mailto:g.fiorenza1@campus.unimib.it}{\texttt{g.fiorenza1@campus.unimib.it}}\\
	\textbf{Ghislotti Gianluca} & 859242 & \href{mailto:g.ghislotti@campus.unimib.it}{\texttt{g.ghislotti@campus.unimib.it}}\\
\end{tabular}

\vfill
\begin{center}
	\large{\textbf{Anno Accademico 2022/2023}}
\end{center}




	\thispagestyle{empty}
	\newpage
	\pagestyle{fancy}
	\pagenumbering{roman}
	
	
	\tableofcontents
	
	\newpage
	
	\definecolor{fireenginered}{rgb}{0.81, 0.09, 0.13}
	\hypersetup{
		linkbordercolor=fireenginered
	}
	\pagenumbering{arabic}
	
	
	\part{Implementazione della DCT2}
	
	Lo scopo di questa parte del progetto è stato quello di implementare la nostra versione della DCT2 e confrontare i tempi di esecuzione con quelli di una libreria che implementa la versione fast della DCT2. La libreria in questione si chiama FFTW \cite{fftw} ed è scritta in C++.
	
	
	
	\part{Progettazione di un'interfaccia grafica per la compressione di immagini}
	\FloatBarrier
	\newpage
	\printbibliography
\end{document}
	
